\documentclass[12pt,a4paper]{article}

% ---- Encoding & Font ----
\usepackage{iftex}
\ifPDFTeX
  \usepackage[utf8]{inputenc}
  \usepackage[T1]{fontenc}
\else
  \usepackage{fontspec}  % XeTeX/LuaTeX (tectonic): Unicode nativo (·, —, ¿, ª, §)
\fi
\usepackage[spanish]{babel}
\usepackage{csquotes}

% ---- Page layout ----
\usepackage[top=2.5cm, bottom=2.5cm, left=2.5cm, right=2.5cm]{geometry}
\usepackage{setspace}
\onehalfspacing

% ---- Math ----
\usepackage{amsmath, amssymb, amsthm}
\usepackage{mathtools}
\usepackage{physics}
\usepackage{esint}  % \oiint

% ---- Graphics ----
\usepackage{graphicx}
\usepackage{tikz}
\usepackage{float}
\usepackage{caption}

% ---- Tables ----
\usepackage{array}
\usepackage{booktabs}
\usepackage{tabularx}
\usepackage{colortbl}

% ---- Colours ----
\usepackage{xcolor}
\definecolor{notebg}{HTML}{FFF3CD}
\definecolor{noteborder}{HTML}{FFC107}
\definecolor{boxred}{HTML}{F8D7DA}
\definecolor{boxredborder}{HTML}{F5C6CB}

% ---- Boxes ----
\usepackage{mdframed}
\usepackage{tcolorbox}
\tcbuselibrary{most}

\newtcolorbox{keybox}[1][]{
  colback=blue!5,
  colframe=blue!70!black,
  arc=4pt,
  boxrule=1pt,
  fonttitle=\bfseries,
  title={#1}
}

\newtcolorbox{warnbox}[1][]{
  colback=boxred,
  colframe=boxredborder,
  coltitle=black!75,
  arc=4pt,
  boxrule=1pt,
  fonttitle=\bfseries,
  title={#1}
}

\newtcolorbox{notebox}[1][]{
  colback=notebg,
  colframe=noteborder,
  coltitle=black!75,
  arc=4pt,
  boxrule=1pt,
  fonttitle=\bfseries,
  title={#1}
}

% ---- Hyperlinks & TOC ----
\usepackage[hidelinks]{hyperref}
\usepackage{bookmark}

% ---- Enumerate ----
\usepackage{enumitem}

% ---- Miscellaneous ----
\usepackage{icomma}
\usepackage{siunitx}
\usepackage{textcomp}
\usepackage[bottom]{footmisc}

% ---- Title ----
\title{\textbf{Clase 19: Varilla Rotante, Generadores y Motores, \\
        Campos Eléctricos Inducidos e Inductancia}\\[0.3em]
        \large{FEM de movimiento, corriente alterna, Faraday para $\mathbf{E}$
              y autoinductancia}}
\author{Transcripto y expandido --- Curso Física III (OpenFING)}
\date{}

\begin{document}

\maketitle
\thispagestyle{empty}
\newpage

\tableofcontents
\newpage

% ===================================================================
\section{Varilla que rota en un campo magnético}
% ===================================================================

Varilla conductora de largo $R$, con un extremo fijo, que rota con velocidad
angular $\omega$ en un campo uniforme $B$ perpendicular. Se busca la FEM entre
sus extremos, por dos métodos.

\subsection{Método 1: suma de FEMs de segmentos}

No se puede usar $\varepsilon = LvB$ directamente porque la velocidad depende
del punto ($v = \omega r$). Un segmento $dr$ a distancia $r$ aporta
$d\varepsilon = B(\omega r)\,dr$. En serie, se integran:
\[
\varepsilon = \int_0^{R} B\,\omega\, r\, dr = \frac{B\,\omega\,R^2}{2}.
\]

\begin{keybox}[FEM de la varilla rotante]
  \[
  \boxed{|\varepsilon| = \frac{B\,\omega\,R^2}{2}}
  \]
\end{keybox}

\subsection{Método 2: Faraday sobre un circuito}

Cerrando el circuito, el área barrida es $A = \tfrac{\theta}{2}R^2$ con
$\theta = \omega t$, así que $\Phi_B = \tfrac{B\omega R^2}{2}t$ y
\[
\left|\frac{d\Phi_B}{dt}\right| = \frac{B\,\omega\,R^2}{2},
\]
idéntico al método 1. Coherente.

% ===================================================================
\section{Generadores y motores}
% ===================================================================

\subsection{El generador de corriente alterna}

Campo uniforme $B$ y espira (bobina de $N$ vueltas) girando con $\omega$
(turbina). Con $\theta = \omega t$ el ángulo normal--$\mathbf{B}$:
\[
\Phi_B = B\,A\,\cos(\omega t).
\]

\begin{keybox}[FEM sinusoidal del generador]
  \[
  \boxed{\varepsilon = -\frac{d\Phi_B}{dt} = B\,A\,\omega\,\sin(\omega t)}
  \]
\end{keybox}

Conectada a $R$, $I = \tfrac{BA\omega}{R}\sin(\omega t)$: \textbf{corriente
alterna}. Los generadores simples producen CA de forma natural.

\subsection{Torque necesario y balance de energía}

Para $\omega$ constante hace falta un torque externo opuesto al magnético
($\boldsymbol\tau = I\mathbf{A}\times\mathbf{B}$):
\[
\tau = I A B \sin\theta = \frac{B^2 A^2 \omega}{R}\sin^2\theta.
\]
El generador convierte energía mecánica en eléctrica.

\subsection{El motor: el mismo aparato al revés}

Inyectando una corriente alterna, el torque hace girar el eje: es un
\textbf{motor} (energía eléctrica $\to$ mecánica).

\begin{notebox}[¿Por qué CA?]
  Con corriente continua el torque $\propto\sin\theta$ cambia de signo cada
  media vuelta y no hay efecto acumulado. La CA sincroniza el cambio de signo;
  el motor DC usa un conmutador.
\end{notebox}

% ===================================================================
\section{Campos eléctricos inducidos}
% ===================================================================

El campo magnético no trabaja: para mover cargas hace falta un campo eléctrico
inducido asociado a la FEM.

\subsection{Faraday en términos del campo eléctrico}

Una carga que recorre $C$ recibe $\varepsilon = \oint_C \mathbf{E}\cdot
d\mathbf{l}$. Igualando a Faraday:

\begin{keybox}[Ley de Faraday para $\mathbf{E}$]
  \[
  \boxed{\oint_C \mathbf{E}\cdot d\mathbf{l} = -\frac{d\Phi_B}{dt}}
  \]
\end{keybox}

\begin{warnbox}[$\mathbf{E}$ ya no deriva de un potencial]
  Con campo magnético variable, la circulación de $\mathbf{E}$ no es cero:
  $\mathbf{E} = -\nabla V$ deja de valer. Siguen valiendo $\mathbf{F}=q_0
  \mathbf{E}$, el trabajo y la ley de Ohm.
\end{warnbox}

\subsection{Ejemplo: campo B uniforme y variable}

Región circular con $\mathbf{B}$ uniforme pero variable, sin cargas. Las líneas
de $\mathbf{E}$ son cerradas (círculos concéntricos), $\mathbf{E}$ tangente y de
módulo constante. Con $C$ de radio $r$:
\[
E\,(2\pi r) = -\frac{d\Phi_B}{dt}, \qquad \Phi_B = B\,\pi r^2.
\]

\begin{keybox}[Campo eléctrico inducido]
  \[
  \boxed{|E| = \frac{r}{2}\left|\frac{dB}{dt}\right|}
  \]
\end{keybox}

Crece linealmente con $r$; los campos inducidos no derivan de un potencial.

% ===================================================================
\section{Inductancia}
% ===================================================================

La inductancia es a las bobinas lo que la capacitancia a los condensadores:
caracteriza la energía almacenable en el campo magnético.

\subsection{Definición y autoinductancia}

Se llama \textbf{autoinductancia}: una bobina se induce FEM a sí misma al
variar su corriente (distinta de la inductancia mutua).

\subsection{Cálculo para un solenoide ideal}

Solenoide ideal ($n$ espiras/longitud, $N = n L_{\text{sol}}$, sección $A$),
$B = \mu_0 n I$. Flujo en una espira $\Phi_1 = \mu_0 n I A$; FEM total:
\[
|\varepsilon| = \mu_0\, n^2\, A\, L_{\text{sol}}\,\left|\frac{dI}{dt}\right|.
\]

\begin{keybox}[Autoinductancia de un solenoide]
  \[
  \boxed{L = \mu_0\, n^2\, A\, L_{\text{sol}} = \frac{\mu_0 N^2 A}{L_{\text{sol}}}}
  \]
\end{keybox}

\subsection{Caso general, unidades y signos}

Como $\Phi_B \propto I$ y la bobina es rígida, $\Phi_B = L I$ y
\[
\varepsilon = -L\,\frac{dI}{dt}.
\]
Receta: $L = \Phi_B^{\text{total}}/I$. Unidad: el \textbf{henry} (H) $=$
V$\cdot$s/A.

\begin{keybox}[Signos en una bobina]
  \[
  \boxed{V_B - V_A = -L\,\frac{dI}{dt}}
  \]
  El signo es el mismo crezca o decrezca $I$ (análogo a $-RI$).
\end{keybox}

\subsection{Ejemplo: el toroide}

Toroide de sección rectangular (altura $h$, radios $A$ y $B$, $N$ espiras),
$B(r) = \mu_0 N I/(2\pi r)$. Flujo en una espira:
\[
\Phi_1 = \frac{h N \mu_0 I}{2\pi}\ln\!\frac{B}{A}.
\]

\begin{keybox}[Autoinductancia de un toroide]
  \[
  \boxed{L = \frac{\mu_0 N^2 h}{2\pi}\ln\!\frac{B}{A}}
  \]
\end{keybox}

\begin{notebox}[Próxima clase]
  Circuitos RL, materiales magnéticos y energía en el campo magnético.
\end{notebox}

\end{document}
